% Authored: how the manual is organised and how it is produced.

\section*{The three questions}

Technical documentation of a repository has to answer three distinct
questions, and mixing them produces a document that answers none of them well.
This manual keeps them apart:

\begin{description}
  \item[Reference Manual] \emph{What does the repository expose?} An
    enumeration of every public object with its exact signature, its
    parameters, its failure modes, and the source file that defines it.
  \item[Software Architecture] \emph{How is the system structured, and how do
    its parts interact?} Components, dependency direction, runtime flow, and
    data lineage.
  \item[Engineering] \emph{How is the system developed, tested, built,
    reproduced, packaged, and released?} Structure, workflow, quality gates,
    CI/CD, reproducibility, and releases.
\end{description}

\begin{figure}[htbp]
\centering
\begin{tikzpicture}[
  every node/.style={font=\small},
  box/.style={rectangle, draw=black, rounded corners=2pt, align=center,
    minimum width=30mm, minimum height=8mm},
  leaf/.style={rectangle, draw=rulegrey, align=center, font=\scriptsize,
    minimum width=24mm, minimum height=6mm},
  link/.style={-{Latex[length=2mm]}, thick},
  thin link/.style={-{Latex[length=1.6mm]}, rulegrey}
]
  \node[box] (root) {Project documentation};
  \node[box, below=16mm of root] (arch) {Architecture};
  \node[box, left=14mm of arch] (ref) {Reference manual};
  \node[box, right=14mm of arch] (eng) {Engineering};

  \node[leaf, below=8mm of ref] (r1) {Modules, classes,\\ functions};
  \node[leaf, below=5mm of r1] (r2) {Schemas, CLI, HTTP,\\ configuration};
  \node[leaf, below=8mm of arch] (a1) {Components,\\ dependencies};
  \node[leaf, below=5mm of a1] (a2) {Runtime flow,\\ data lineage};
  \node[leaf, below=8mm of eng] (e1) {Tests, quality gates,\\ CI/CD};
  \node[leaf, below=5mm of e1] (e2) {Reproducibility,\\ releases};

  \node[box, below=12mm of a2, minimum width=46mm] (src) {Source repository\\
    \scriptsize (single source of truth)};

  \draw[link] (root) -- (ref);
  \draw[link] (root) -- (arch);
  \draw[link] (root) -- (eng);
  \draw[thin link] (ref) -- (r1);
  \draw[thin link] (r1) -- (r2);
  \draw[thin link] (arch) -- (a1);
  \draw[thin link] (a1) -- (a2);
  \draw[thin link] (eng) -- (e1);
  \draw[thin link] (e1) -- (e2);
  \draw[link] (src) -- (a2);
  \draw[link] (src.west) to[out=170, in=-90] (r2.south);
  \draw[link] (src.east) to[out=10, in=-90] (e2.south);
\end{tikzpicture}
\caption{Information architecture of this manual. All three references derive
from the same source repository; none of them is maintained by hand where the
information can be extracted.}
\label{fig:information-architecture}
\end{figure}

\section*{Authored and generated content}

Two kinds of content appear in this manual and they are never mixed in the same
file.

\emph{Generated} content is derived from the repository by the tools in
\srcfile{docs/tools/} and written under \srcfile{docs/latex/generated/}. Every
generated file carries a header marking it as machine-produced. Signatures,
parameter tables, field tables, CLI options, endpoint tables, environment
variables, dependency lists, the module graph, the file census, and all indexes
are generated. Editing them by hand is pointless: the next \texttt{make docs}
overwrites them.

\emph{Authored} content is the prose you are reading: chapter introductions,
the mathematical exposition in \cref{ch:mathematics}, the architecture
narrative, the architecture decision records, and the engineering commentary.
Authored files live directly under \srcfile{docs/latex/reference/},
\srcfile{docs/latex/architecture/}, \srcfile{docs/latex/engineering/}, and
\srcfile{docs/latex/diagrams/}, and the generator never writes to them.

\section*{Two documents}

The documentation is compiled into two PDFs from one source tree. They share
this preamble, the generated metadata, and every chapter file, so they cannot
describe different revisions and nothing is written twice.

\begin{description}[nosep]
  \item[\srcfile{docs/latex/reference.tex}] The documented objects only: the
    CRAN-style reference, for a reader who wants to look an object up.
  \item[\srcfile{docs/latex/main.tex}] This document: the same reference, plus
    the software architecture and engineering parts and the traceability
    appendix.
\end{description}

Entries use the layout CRAN reference manuals use: the object name as the
heading, its summary immediately below, then run-in field labels
(\textbf{Usage}, \textbf{Arguments}, \textbf{Value}, \textbf{Raises},
\textbf{Details}, \textbf{Source}). Arguments are hanging-indent lists rather
than full-width tables, which is what keeps a reference of this size readable
at a sane page count.

\section*{Regenerating this manual}

From the repository root:

\begin{lstlisting}[language=bash]
make docs                # inspect, extract, render, validate
make docs-pdf            # the above, then build the complete manual
make docs-pdf-reference  # build the reference-only document
\end{lstlisting}

\texttt{make docs} performs no application imports. It reads the source with
Python's \texttt{ast} module, so generating documentation cannot open a network
connection, contact an LLM provider, build a vector index, or write to any
external system.

\section*{Determinism and revision correspondence}

Two runs of the generator over the same repository revision produce
byte-identical output. Timestamps come from the HEAD commit date rather than
the wall clock, object order is by fully-qualified name, and every table row
order is derived from the model rather than from dictionary iteration.

The manual therefore describes exactly one revision, recorded on the title page
and in \srcfile{docs/metadata/documentation-manifest.json} together with a
SHA-256 digest of every generated file. If the source has moved on, the digests
no longer match and \texttt{python -m tools.validate\_docs} reports the drift.

\section*{Secrets}

No secret value appears in this manual. The configuration reference documents
the \emph{names}, types, defaults, and sensitivity of credential settings; the
declared default of every credential field in source is \texttt{None}, and the
generator reads source text rather than the process environment, so no
environment value can reach the output. The validator additionally scans every
generated file for credential-shaped literals.
